% Created 2025-10-23 jue 08:58
% Intended LaTeX compiler: pdflatex
\documentclass[a4paper]{article}
\usepackage[utf8]{inputenc}
\usepackage[T1]{fontenc}
\usepackage{graphicx}
\usepackage{longtable}
\usepackage{wrapfig}
\usepackage{rotating}
\usepackage[normalem]{ulem}
\usepackage{amsmath}
\usepackage{amssymb}
\usepackage{capt-of}
\usepackage{hyperref}
\usepackage[spanish]{babel}
\usepackage[usenames,dvipsnames]{color} % Required for custom colors
\renewcommand{\ttdefault}{pcr} % MONOESPACIO CON NEGRIT
\usepackage{lastpage}
\usepackage{listings}
\usepackage{listingsutf8}
\renewcommand{\lstlistingname}{Listado}
\lstset{frame=single,inputencoding=utf8,basicstyle=\scriptsize\ttfamily,showstringspaces=false,numbers=none}
\definecolor{MyDarkGreen}{rgb}{0.0,0.4,0.0} % This is the color used for comments
\lstset{ breaklines=true, postbreak=\mbox{\textcolor{red}{$\hookrightarrow$}\space}, keywordstyle=\bfseries, keywordstyle=[1]\color{Blue}\bfseries,  keywordstyle=[2]\color{Purple}\bfseries,  keywordstyle=[3]\color{Blue}\underbar,   identifierstyle=,   commentstyle=\usefont{T1}{pcr}{m}{sl}\color{MyDarkGreen}\small,   stringstyle=\color{Purple},   showstringspaces=false,   tabsize=2,   morecomment=[l][\color{Blue}]{...} }
\lstset{literate=  {á}{{\'a}}1 {é}{{\'e}}1 {í}{{\'i}}1 {ó}{{\'o}}1 {ú}{{\'u}}1   {Á}{{\'A}}1 {É}{{\'E}}1 {Í}{{\'I}}1 {Ó}{{\'O}}1 {Ú}{{\'U}}1   {à}{{\`a}}1 {è}{{\`e}}1 {ì}{{\`i}}1 {ò}{{\`o}}1 {ù}{{\`u}}1   {À}{{\`A}}1 {È}{{\'E}}1 {Ì}{{\`I}}1 {Ò}{{\`O}}1 {Ù}{{\`U}}1   {ä}{{\"a}}1 {ë}{{\"e}}1 {ï}{{\"i}}1 {ö}{{\"o}}1 {ü}{{\"u}}1   {Ä}{{\"A}}1 {Ë}{{\"E}}1 {Ï}{{\"I}}1 {Ö}{{\"O}}1 {Ü}{{\"U}}1   {â}{{\^a}}1 {ê}{{\^e}}1 {î}{{\^i}}1 {ô}{{\^o}}1 {û}{{\^u}}1   {Â}{{\^A}}1 {Ê}{{\^E}}1 {Î}{{\^I}}1 {Ô}{{\^O}}1 {Û}{{\^U}}1   {œ}{{\oe}}1 {Œ}{{\OE}}1 {æ}{{\ae}}1 {Æ}{{\AE}}1 {ß}{{\ss}}1   {ű}{{\H{u}}}1 {Ű}{{\H{U}}}1 {ő}{{\H{o}}}1 {Ő}{{\H{O}}}1   {ç}{{\c c}}1 {Ç}{{\c C}}1 {ø}{{\o}}1 {å}{{\r a}}1 {Å}{{\r A}}1   {€}{{\euro}}1 {£}{{\pounds}}1 {«}{{\guillemotleft}}1   {»}{{\guillemotright}}1 {ñ}{{\~n}}1 {Ñ}{{\~N}}1 {¿}{{?`}}1 }
\usepackage{caption}
\usepackage{attachfile2}
\usepackage[margin=1.5cm,includeheadfoot,includehead,includefoot]{geometry}
\hypersetup{colorlinks,linkcolor=black}
\usepackage{fancyhdr}
\pagestyle{fancyplain}
\chead{}
\lhead{}
\rhead{}
\cfoot{}
\lfoot{\begin{footnotesize}alvaro.gonzalezsotillo@educa.madrid.org\end{footnotesize}}
\rfoot{\begin{footnotesize}\thepage / \pageref{LastPage}\end{footnotesize}}
\usepackage{svg}
\usepackage{letltxmacro}
\LetLtxMacro{\originalincludegraphics}{\includegraphics}
\renewcommand{\includegraphics}[2][]{\IfFileExists{#2.pdf}{\originalincludegraphics[#1]{#2.pdf}}{\originalincludegraphics[#1]{#2}}}
\LetLtxMacro{\originalincludesvg}{\includesvg}
\renewcommand{\includesvg}[2][]{\IfFileExists{#2.pdf}{\originalincludegraphics[#1]{#2.pdf}}{\originalincludegraphics[#1]{#2.svg.pdf}}}
\usepackage{comment}
\excludecomment{NOTES}
\author{Álvaro González Sotillo}
\date{\today}
\title{Elementos del desarrollo de Software}
\hypersetup{
 pdfauthor={Álvaro González Sotillo},
 pdftitle={Elementos del desarrollo de Software},
 pdfkeywords={},
 pdfsubject={},
 pdfcreator={Emacs 30.2 (Org mode 9.7.11)}, 
 pdflang={Spanish}}
\begin{document}

\maketitle
\setcounter{tocdepth}{1}
\tableofcontents

\captionsetup{font=scriptsize}
\section{Introducción}
\label{sec:org000000c}
\begin{itemize}
\item En esta Unidad aprenderemos a
\begin{itemize}
\item Reconocer la relación de los programas con los componentes del sistema
informático.
\item Diferenciar código fuente, objeto y ejecutable.
\item Identificar las fases de desarrollo de una aplicación informática.
\item Clasificar los lenguajes de programación.
\end{itemize}
\end{itemize}
\subsection{Tipos de software}
\label{sec:org0000000}
\begin{itemize}
\item \textbf{De sistema} (Sistema operativo, \emph{drivers} -controladores-)
\begin{itemize}
\item Los \textbf{\emph{drivers}} son los controladores de dispositivos.
\item El \textbf{\emph{firmware}} es el programa que funciona dentro de un dispositivo
\item \textbf{BIOS/UEFI} es el \emph{firmware} de la placa base de un PC
\end{itemize}
\item \textbf{De aplicación} (Suite ofimática, Navegador, Edición de imagen, \ldots{})
\item \textbf{De desarrollo} (Editores, compiladores, intérpretes, \ldots{})
\end{itemize}

\begin{center}
\includegraphics[width=.9\linewidth]{/home/alvaro/.cache/org-persist/1d/e43e17-23e7-4760-bf76-d8155651f27f-a4d0a1f3ca59e90f59626b9721b0aac2.png}
\end{center}
\subsection{Relación Hardware-Software}
\label{relación-hardware-software}
\begin{itemize}
\item \textbf{Disco duro} : almacena de forma permanente los archivos ejecutables y los archivos de datos.
\begin{itemize}
\item El disco duro se considera un periférico de E/S (Entrada/Salida).
\end{itemize}
\item \textbf{Memoria RAM} : almacena de forma temporal el código binario de los archivos ejecutables y los archivos de datos necesarios.
\item \textbf{CPU} : lee y ejecuta instrucciones almacenadas en memoria RAM, así  como los datos necesarios.
\item \textbf{E/S}: recoge nuevos datos desde la entrada, se muestran los resultados, se leen/guardan a disco, \ldots{}
\end{itemize}
\subsection{¿pero qué es un ordenador?}
\label{sec:org0000003}
\begin{itemize}
\item Es un dispositivo capaz de
\begin{itemize}
\item Ejecutar órdenes de forma secuencial
\item Con dos tipos especiales de órdenes
\begin{itemize}
\item Condicionales: se ejecutan dependiendo de los datos disponibles
\item Salto: se interrumpe el orden normal de ejecución
\end{itemize}
\end{itemize}

\item Preguntas:
\begin{itemize}
\item ¿una reloj digital es un ordenador?
\item ¿una \href{https://www.casio.com/es/scientific-calculators/product.FX-55PLUS/}{calculadora} es un \href{https://www.casio.com/es/scientific-calculators/product.FX-9860GIII/}{ordenador}?
\item ¿\href{https://bogdanthegeek.github.io/blog/projects/vapeserver/}{un vaper es un ordenador}?
\end{itemize}
\end{itemize}

\begin{itemize}
\item Conclusión
\begin{itemize}
\item Muchísimos dispositivos contienen ordenadores
\item No todos esos ordenadores se exponen al usuario final
\begin{itemize}
\item El usuario final no puede modificar los programas
\end{itemize}
\end{itemize}
\end{itemize}
\subsection{Arquitectura de Von neuman}
\label{sec:org0000006}
\begin{itemize}
\item La \href{https://es.wikipedia.org/wiki/M\%C3\%A1quina\_de\_Turing\#Descripci\%C3\%B3n\_informal}{máquina de Turing} distinguía entre datos y programa
\item Pero los programas son también datos. Esto posibilita:
\begin{itemize}
\item Emuladores
\item Programas que hacen programas
\item Actualización de software
\end{itemize}
\end{itemize}

\begin{figure}[htbp]
\centering
\includegraphics[width=.9\linewidth]{/home/alvaro/.cache/org-persist/85/d563ae-624c-44d8-acdc-4ad2049a67eb-2f212604807be4c9ebe3316d8ea9f126.png}
\caption{Arquitectura de Von Neuman}
\end{figure}
\subsection{Códigos fuente, objeto y ejecutable}
\label{códigos-fuente-objeto-y-ejecutable}
\begin{itemize}
\item \textbf{Código fuente}: archivo de texto legible escrito en un lenguaje de programación.
\item \textbf{Código objeto} (intermedio): archivo binario no ejecutable.
\begin{itemize}
\item Pueden ser ya instrucciones legibles por la CPU
\item O puede ser un formato intermedio
\end{itemize}
\item \textbf{Código ejecutable}: código máquina
\begin{itemize}
\item En lenguajes interpretados no existe código objeto, ni ejecutable. Solo código fuente. Por ejemplo: PHP, Javascript, \ldots{}
\item Aunque podría existir en memoria (\href{https://en.wikipedia.org/wiki/Just-in-time\_compilation}{JIT compiler})
\end{itemize}
\end{itemize}

\begin{lstlisting}[language=asm,caption={Código ensamblador amd64, y su equivalente a código máquina},captionpos=b,numbers=none]
objdump -d  /usr/bin/ls


/usr/bin/ls:     file format elf64-x86-64


Disassembly of section .init:

000000000000033c <.init>:
 33c:   f3 0f 1e fa             endbr64
 340:   48 83 ec 08             sub    $0x8,%rsp
 344:   48 8b 05 75 2c 02 00    mov    0x22c75(%rip),%rax        # 22fc0 <__gmon_start__@Base>
 34b:   48 85 c0                test   %rax,%rax
 34e:   74 02                   je     352 <__ctype_toupper_loc@plt-0x67e>
 350:   ff d0                   call   *%rax
 352:   48 83 c4 08             add    $0x8,%rsp
 356:   c3                      ret

Disassembly of section .plt:

0000000000000360 <.plt>:
 360:   ff 35 ca 28 02 00       push   0x228ca(%rip)        # 22c30 <error_at_line@@Base+0x21943>
 366:   ff 25 cc 28 02 00       jmp    *0x228cc(%rip)        # 22c38 <error_at_line@@Base+0x2194b>
 36c:   0f 1f 40 00             nopl   0x0(%rax)
 370:   f3 0f 1e fa             endbr64
 374:   68 00 00 00 00          push   $0x0
 379:   e9 e2 ff ff ff          jmp    360 <__ctype_toupper_loc@plt-0x670>
 37e:   66 90                   xchg   %ax,%ax
 380:   f3 0f 1e fa             endbr64
        
 ...  
\end{lstlisting}
\subsection{Para entenderlo mejor}
\label{sec:org0000009}

\begin{itemize}
\item Simulador de CPU: \url{https://cpuvisualsimulator.github.io/}
\item Compilador online: \href{https://godbolt.org/}{https://godbolt.org/}
\end{itemize}
\section{Ciclo de vida del software}
\label{ciclo-de-vida-del-software}
\subsection{Ingeniería de software}
\label{ingeniería-de-software}
\begin{quote}
Disciplina que estudia los principios y metodologías para el desarrollo y mantenimiento de sistemas software.
\end{quote}

\begin{itemize}
\item Algunos autores consideran que "\textbf{\emph{desarrollo de software}}" es un término más apropiado que "\textbf{\emph{ingeniería de software}}" puesto que este último implica niveles de rigor y prueba de procesos que no son apropiados para todo tipo de desarrollo de software.
\item \href{https://es.wikipedia.org/wiki/Ingenier\%C3\%ADa\_de\_software}{Ingeniería del software}
\end{itemize}
\subsection{Desarrollo de software}
\label{desarrollo-de-software}
\subsubsection{\textbf{\emph{Fases principales}}}
\label{fases-principales}
\begin{itemize}
\item \textbf{ANÁLISIS}: Qué debe hacer la aplicación.
\item \textbf{DISEÑO}: Cómo se va a organizar internamente la aplicación. Sus decisiones no afectan directamente al usario, más allá de elegir una plataforma.
\item \textbf{CODIFICACIÓN}: Realización de la aplicación.
\item \textbf{PRUEBAS}: Búsqueda de defectos.
\item \textbf{MANTENIMIENTO}: Solución de defectos que escaparon a las pruebas, o adaptación a nuevos requisitos.
\end{itemize}
\subsubsection{ANÁLISIS}
\label{análisis}
\begin{itemize}
\item Se determina y define claramente las necesidades del cliente y se
especifica los requisitos que debe cumplir el software a desarrollar.
\item La \textbf{especificación de requisitos} debe:
\begin{itemize}
\item Ser completa y sin omisiones
\item Ser concisa y sin trivialidades
\item Evitar ambigüedades. Utilizar lenguaje formal.
\item Evitar detalles de diseño o implementación
\item Ser entendible por el cliente
\item Separar requisitos funcionales y no funcionales
\item Dividir y jerarquizar el modelo
\item Fijar criterios de validación
\end{itemize}
\end{itemize}
\subsubsection{DISEÑO}
\label{diseño}
\begin{itemize}
\item Se descompone y organiza el sistema en elementos componentes que
pueden ser desarrollados por separado.
\item Se especifica la interrelación y funcionalidad de los elementos
componentes.
\item Las actividades habituales son las siguientes:
\begin{itemize}
\item Diseño arquitectónico
\item Diseño detallado
\item Diseño de datos
\item Diseño de interfaz de usuario
\end{itemize}
\end{itemize}
\subsubsection{CODIFICACIÓN}
\label{codificación}
\begin{itemize}
\item Se escribe el código fuente de cada componente.
\item Pueden utilizarse distintos lenguajes informáticos y plataformas:
\begin{itemize}
\item Aplicación web: javascript, HTML, PHP, JSP, ASP\ldots{}
\item Escritorio: C++, python, java, C\#
\item Móvil: flutter, kotlin\ldots{}
\end{itemize}
\end{itemize}
\subsubsection{PRUEBAS}
\label{pruebas}
\begin{itemize}
\item El principal objetivo de las pruebas debe ser conseguir que el
programa funcione incorrectamente y que se descubran defectos.
\item Deberemos someter al programa al máximo número de situaciones
diferentes.
\end{itemize}
\subsubsection{MANTENIMIENTO}
\label{mantenimiento}
\begin{itemize}
\item Durante la explotación del sistema software es necesario realizar
cambios ocasionales.
\item Para ello hay que rehacer parte del trabajo realizado en las fases
previas.
\item Tipos de mantenimiento:
\begin{itemize}
\item \textbf{Correctivo}: se corrigen defectos.
\item \textbf{Perfectivo}: se mejora la funcionalidad.
\item \textbf{Evolutivo}: se añade funcionalidades nuevas.
\item \textbf{Adaptativo}: se adapta a nuevos entornos.
\end{itemize}
\end{itemize}

¿Qué tipos de mantenimiento suponen una nueva fase de análisis y diseño?    
\subsection{Resultado tras cada fase (I)}
\label{resultado-tras-cada-fase-i}
\begin{itemize}
\item Ingeniería de sistemas: \textbf{Especificación del sistema}
\item ANÁLISIS: \textbf{Especificación de requisitos del software}
\item DISEÑO
\begin{itemize}
\item arquitectónico: \textbf{Documento de arquitectura del software}
\item detallado: \textbf{Especificación de módulos y funciones}
\end{itemize}
\item CODIFICACIÓN: \textbf{Código fuente}
\end{itemize}
\subsection{Resultado tras cada fase (II)}
\label{resultado-tras-cada-fase-ii}
( \emph{Continuación} )

\begin{itemize}
\item PRUEBAS de unidades: \textbf{Módulos utilizables}
\item PRUEBAS de integración: \textbf{Sistema utilizable}
\item PRUEBAS del sistema: \textbf{Sistema aceptado}
\item Documentación: \textbf{Documentación técnica y de usuario}
\item MANTENIMIENTO: \textbf{Informes de errores y control de cambios}
\end{itemize}
\subsection{Modelos de desarrollo de software}
\label{modelos-de-desarrollo-de-software}
\begin{itemize}
\item Modelos clásicos (predictivos)
\begin{itemize}
\item Modelo en cascada
\item Modelo en V
\end{itemize}
\item Modelo de construcción de prototipos
\item Modelos evolutivos o incrementales
\begin{itemize}
\item Modelo en espiral (iterativos)
\item Metodologías ágiles (adaptativos)
\end{itemize}
\end{itemize}
\subsection{Modelo en cascada (I)}
\label{modelo-en-cascada-i}
\url{http://jamj2000.github.io/entornosdesarrollo/1/assets/cascada.png}
\subsection{Modelo en cascada (II)}
\label{modelo-en-cascada-ii}
\begin{itemize}
\item Modelo de mayor antigüedad.
\item Identifica las fases principales del desarrollo software.
\item Las fases han de realizarse en el orden indicado.
\item El resultado de una fase es la entrada de la siguiente fase.
\item Es un modelo bastante rígido que se adapta mal al cambio continuo de
especificaciones.
\item Existen diferentes variantes con mayor o menor cantidad de
actividades.
\end{itemize}
\subsection{Modelo en V (I)}
\label{modelo-en-v-i}
\url{http://jamj2000.github.io/entornosdesarrollo/1/assets/v.png}
\subsection{Modelo en V (II)}
\label{modelo-en-v-ii}
\begin{itemize}
\item Modelo muy parecido al modelo en cascada.
\item Visión jerarquizada con distintos niveles.
\item Los niveles superiores indican mayor abstracción.
\item Los niveles inferiores indican mayor nivel de detalle.
\item El resultado de una fase es la entrada de la siguiente fase.
\item Existen diferentes variantes con mayor o menor cantidad de
actividades.
\end{itemize}
\subsection{Prototipos (I)}
\label{prototipos-i}
\begin{itemize}
\item A menudo los requisitos no están especificados claramente:
\begin{itemize}
\item por no existir experiencia previa.
\item por omisión o falta de concreción del usuario/cliente.
\end{itemize}
\end{itemize}

\url{http://jamj2000.github.io/entornosdesarrollo/1/assets/prototipos.png}
\subsection{Prototipos (II)}
\label{prototipos-ii}
\begin{itemize}
\item Proceso:
\begin{itemize}
\item Se crea un prototipo durante la \textbf{fase de análisis} y es probado por
el usuario/cliente para refinar los requisitos del software a
desarrollar.
\item Se repite el paso anterior las veces necesarias.
\end{itemize}
\end{itemize}
\subsection{Prototipos (III)}
\label{prototipos-iii}
\begin{itemize}
\item Tipos de prototipos:
\begin{itemize}
\item \textbf{Prototipos rápidos} (desechables)
\begin{itemize}
\item El prototipo puede estar desarrollado usando otro lenguaje y/o
herramientas.
\item Finalmente el prototipo se desecha.
\end{itemize}
\item \textbf{Prototipos evolutivos}
\begin{itemize}
\item El prototipo está diseñado en el mismo lenguaje y herramientas del
proyecto.
\item El prototipo se usa como base para desarrollar el proyecto.
\end{itemize}
\end{itemize}
\end{itemize}
\subsection{Modelo en espiral (I)}
\label{modelo-en-espiral-i}
\begin{itemize}
\item Desarrollado por Boehm en 1988.
\item La actividad de \textbf{\emph{ingeniería}} corresponde a las fases de los modelos
clásicos: análisis, diseño, codificación, \ldots{}
\end{itemize}

\url{http://jamj2000.github.io/entornosdesarrollo/1/assets/espiral.png}
\subsection{Modelo en espiral (II)}
\label{modelo-en-espiral-ii}
\begin{itemize}
\item En la actividad de \textbf{\emph{ingeniería}} se da gran importancia a la
reutilización de código.
\end{itemize}

\url{http://jamj2000.github.io/entornosdesarrollo/1/assets/espiral-poo.png}
\subsection{Metodologías ágiles (I)}
\label{metodologías-ágiles-i}
\begin{itemize}
\item Son métodos de ingeniería del software basados en el desarrollo
iterativo e incremental.
\item Los requisitos y soluciones evolucionan con el tiempo según la
necesidad del proyecto.
\item El trabajo es realizado mediante la colaboración de equipos
auto-organizados y multidisciplinarios, inmersos en un proceso
compartido de toma de decisiones a corto plazo.
\item Las metodologías más conocidas son:
\begin{itemize}
\item Kanban
\item Scrum
\item XP (eXtreme Programming)
\end{itemize}
\end{itemize}
\subsection{Metodologías ágiles (II)}
\label{metodologías-ágiles-ii}
\href{https://es.wikipedia.org/wiki/Manifiesto\_\%C3\%A1gil}{Manifiesto por el
Desarrollo Ágil}

\begin{itemize}
\item \textbf{Individuos e interacciones} sobre procesos y herramientas
\item \textbf{Software funcionando} sobre documentación extensiva
\item \textbf{Colaboración con el cliente} sobre negociación contractual
\item \textbf{Respuesta ante el cambio} sobre seguir un plan
\end{itemize}
\subsubsection{\textbf{Kanban (I)}}
\label{kanban-i}
\begin{itemize}
\item También se denomina "sistema de tarjetas".
\item \textbf{Desarrollado inicialmente por Toyota para la industria de fabricación
de productos}.
\item \textbf{Controla por demanda} la fabricación de los productos necesarios en
la cantidad y tiempo necesarios.
\item Enfocado a entregar el máximo valor para los clientes, utilizando los
recursos justos.
\item \href{https://es.wikipedia.org/wiki/Lean\_manufacturing}{Lean
manufacturing}
\item \href{https://dosideas.com/noticias/metodologias/184-como-usar-kanban-en-el-desarrollo-de-software}{Kanban
en desarrollo software}
\end{itemize}
\subsubsection{\textbf{Kanban (II)}}
\label{kanban-ii}
Pizarra kanban

\url{http://jamj2000.github.io/entornosdesarrollo/1/assets/kanban-board.png}
\subsubsection{\textbf{Scrum (I)}}
\label{scrum-i}
\begin{itemize}
\item Modelo de desarrollo incremental.
\item Iteraciones (\textbf{sprint}) regulares cada 2 a 4 semanas.
\item Al principio de cada iteración se establecen sus \textbf{objetivos
priorizados} (\textbf{sprint backlog}).
\item Al finalizar cada iteración se obtiene una \textbf{entrega parcial utilizable
por el cliente}.
\item Existen reuniones diarias para tratar la marcha del \emph{sprint}.
\end{itemize}
\subsubsection{\textbf{Scrum (II)}}
\label{scrum-ii}
\url{http://jamj2000.github.io/entornosdesarrollo/1/assets/scrum.png}
\subsubsection{\textbf{XP (Programación extrema) (I)}}
\label{xp-programación-extrema-i}
\textbf{Valores}

\begin{itemize}
\item Simplicidad
\item Comunicación
\item Retroalimentación
\item Valentía o coraje
\item Respeto o humildad
\end{itemize}
\subsubsection{\textbf{XP (Programación extrema) (II)}}
\label{xp-programación-extrema-ii}
\textbf{Características}

\begin{itemize}
\item Diseño sencillo
\item Pequeñas mejoras continuas
\item Pruebas y refactorización
\item Integración continua
\item \textbf{Programación por parejas}
\item \textbf{El cliente se integra en el equipo de desarrollo}
\item Propiedad del código compartida
\item Estándares de codificación
\item 40 horas semanales
\end{itemize}
\section{Lenguajes de programación}
\label{lenguajes-de-programación}
\subsection{Evolución}
\label{sec:org000000f}

\begin{itemize}
\item Primera generación: lenguaje máquina.
\item Segunda generación: se crearon los primeros lenguajes ensambladores.
\item Tercera generación: se crean los primeros lenguajes de alto nivel. Ej. C, Pascal, Cobol…
\item Cuarta generación: son los lenguajes capaces de generar código por si solos, son los llamados RAD, con lo cuales se pueden realizar aplicaciones sin ser un experto en el lenguaje.
\begin{itemize}
\item \href{https://www.nocode.tech/article/what-is-no-code}{Nocode}
\item Excel
\end{itemize}
\item Quinta generación: Basados en restricciones declarativas, para IA clásica (prolog)
\item \href{https://en.wikipedia.org/wiki/Vibe\_coding}{Vibe coding}: utilizar sistemas LLM para la programación
\begin{itemize}
\item Nuevos (apenas dos años reales)
\item Prometedores, pero solo para programadores con experiencia previa
\end{itemize}
\end{itemize}
\subsection{Obtención de código ejecutable}
\label{obtención-de-código-ejecutable}
\begin{itemize}
\item Para obtener código binario ejecutable tenemos 2 opciones:
\begin{itemize}
\item \textbf{Compilar}
\item \textbf{Interpretar}
\end{itemize}
\end{itemize}
\subsection{Proceso de compilación/interpretación}
\label{proceso-de-compilacióninterpretación}
\begin{itemize}
\item La compilación/interpretación del código fuente se lleva a cabo en dos
fases:
\begin{enumerate}
\item \textbf{Análisis léxico}
\item \textbf{Análisis sintáctico}
\end{enumerate}
\item Si no existen errores, se genera el código objeto correspondiente.
\item Un código fuente correctamente escrito no significa que funcione según
lo deseado.
\item No se realiza un análisis semántico.
\end{itemize}
\subsection{Lenguajes compilados}
\label{lenguajes-compilados}
\begin{itemize}
\item Ejemplos: C, C++
\item Principal ventaja: Ejecución muy eficiente.
\item Principal desventaja: Es necesario compilar cada vez que el código
fuente es modificado.
\end{itemize}
\subsubsection{Ejemplo compilación}
\label{sec:org0000012}
\begin{itemize}
\item \texttt{gcc -E}: Preprocessor, but don't compile
\item \texttt{gcc -S}: Compile but don't assemble
\item \texttt{gcc -c}: Preprocess, compile, and assemble, but don't link
\end{itemize}

\begin{lstlisting}[language=cpp,numbers=none]
#include <stdio>

void main(){
  std::printf( "Resultado:%d", 5*7 );
}
\end{lstlisting}
\subsection{Lenguajes interpretados}
\label{lenguajes-interpretados}
\begin{itemize}
\item Ejemplos: PHP, Javascript
\item Principal ventaja: El código fuente se interpreta directamente.
\item Principal desventaja: Ejecución menos eficiente.
\end{itemize}
\subsection{JAVA (I)}
\label{java-i}
\begin{itemize}
\item Lenguajes compilado e interpretado.
\item El código fuente Java \textbf{se compila} y se obtiene un código binario
intermedio denominado \textbf{bytecode}.
\item Puede considerarse código objeto pero destinado a la máquina virtual
de Java en lugar de código objeto nativo.
\item Después este \textbf{bytecode} \textbf{se interpreta} para ejecutarlo.
\end{itemize}
\subsection{JAVA (II)}
\label{java-ii}
\begin{itemize}
\item Ventajas:
\begin{itemize}
\item Estructurado y Orientado a Objetos
\item Relativamente fácil de aprender
\item Buena documentación y base de usuarios
\end{itemize}
\item Desventajas:
\begin{itemize}
\item Menos eficiente que los lenguajes compilados (con JIT, solo inicialmente)
\end{itemize}
\end{itemize}
\subsection{Tipos (I)}
\label{tipos-i}
\begin{itemize}
\item Según la forma en la que operan:
\begin{itemize}
\item \textbf{Declarativos}: indicamos el resultado a obtener sin especificar los
pasos.
\item \textbf{Imperativos}: indicamos los pasos a seguir para obtener un
resultado.
\end{itemize}
\end{itemize}
\subsection{Tipos (II)}
\label{tipos-ii}
\begin{itemize}
\item Tipos de \textbf{lenguajes declarativos}:
\begin{itemize}
\item Lógicos: Utilizan reglas. Ej: Prolog
\item Funcionales: Utilizan funciones. Ej: Lisp, Haskell
\item Algebraicos: Utilizan sentencias. Ej: SQL
\end{itemize}
\item Normalmente son lenguajes interpretados.
\end{itemize}
\subsection{Tipos (III)}
\label{tipos-iii}
\begin{itemize}
\item Tipos de \textbf{lenguajes imperativos}:
\begin{itemize}
\item Estructurados: C
\item Orientados a objetos: Java
\item Multiparadigma: C++, Javascript
\end{itemize}
\item Los lenguajes orientados a objetos son también lenguajes
estructurados.
\item Muchos de estos lenguajes son compilados.
\end{itemize}
\subsection{Tipos (IV)}
\label{tipos-iv}
\begin{itemize}
\item Tipos de lenguajes según nivel de abstracción:
\begin{itemize}
\item Bajo nivel: ensamblador
\item Medio nivel: C
\item Alto nivel: C++, Java
\end{itemize}
\end{itemize}
\subsection{Historia}
\label{sec:org0000015}
\href{http://rigaux.org/language-study/diagram-light.png}{Historia simplificada}

\href{http://rigaux.org/language-study/diagram.png}{Historia detallada}
\subsection{Criterios para la selección de un lenguaje}
\label{criterios-para-la-selección-de-un-lenguaje}
\begin{itemize}
\item Campo de aplicación
\item Experiencia previa
\item Herramientas de desarrollo
\item Documentación disponible
\item Base de usuarios
\item Reusabilidad
\item Transportabilidad
\item Plataforma de ejecución
\item Imposición del cliente
\end{itemize}
\section{Referencias}
\label{sec:org0000018}
\phantomsection
\label{}
\begin{itemize}
\item Formatos:
\begin{itemize}
\item \href{./ED-01-elementos-desarrollo-software.reveal.html}{Transparencias}
\item \href{./ED-01-elementos-desarrollo-software.pdf}{PDF}
\item \href{./ED-01-elementos-desarrollo-software.wp.html}{Página web}
\item \href{./ED-01-elementos-desarrollo-software.epub}{EPUB}
\end{itemize}
\item Creado con:
\begin{itemize}
\item \href{https://www.gnu.org/s/emacs/}{Emacs}
\item \href{https://gitlab.com/oer/org-re-reveal}{org-re-reveal}
\item \href{https://www.latex-project.org/}{Latex}
\end{itemize}
\item Alojado en \href{https://alvarogonzalezsotillo.github.io/apuntes-clase}{Github}
\end{itemize}
\end{document}
